We model the dynamic insurance choices of a panel of consumers, beginning with their first interaction with our firm. In the initial period, each consumer chooses a level of coverage ($y$) along with a choice of whether to opt in to the monitoring program $m \in \{0,1\}$. In each subsequent (renewal) period, monitoring is no longer offered, but the consumer again chooses a coverage level ($y$) and whether to stay with the firm or switch to another ($f$).

Consumers are forward looking with standard von Neumann-Morgenstern utility under constant absolute risk aversion with parameter $\gamma$. As such, they make their decisions in each period $t$ with respect to their expected consumption utility in period $t$ and the expected continuation value from the subsequent periods. At the beginning of each period, each consumer observes a menu of insurance plans along with their respective premiums and coverage limits. If an accident occurs over the course of the period, then the consumer expects to pay an out-of-pocket liability expenditure (henceforth ``OOP'') that depends on the severity of the accident and the coverage limit chosen at the beginning of the period. In addition to these direct monetary costs, the consumer may face {choice frictions} against overriding the default mandatory-minimum plan, opting in to monitoring in the initial period, or switching plans or firms in renewal periods.

Putting these factors together, a representative consumer who makes choice $d=({y,m,f})$ in period $t$ obtains the following flow consumption value $h$, conditional on their characteristics $x_{t}$, their previous period choice $d_{t-1}$ and the realization of accident losses in that period $A_{t}$:
\begin{align}\label{eq:flow_consumption}
h(d|x_{t}, {A}_{t},d_{t-1})= & -\underbrace{{p_{t}}(d | x_{t})}_{\text{price}}-\underbrace{\eta_{t}(d | d_{t-1})}_{\text{choice friction}}-\underbrace{e(d|{A}_{t})}_{\text{OOP}}.
\end{align}

For a given set of observables $x_t$,
each firm $f$ offers a price menu $p$, which specifies prices of a common set of plans $Y$, where $y\in Y$ indexes coverage levels increase monotonically with price. Customers at the monitoring firm may also receive a multiplicative discount or surcharge associated with their monitoring choice and performance. The choice friction $\eta_{t}(d | d_{t-1})$ is composed of several components. In the first period, it includes an additive disutility factor associated with monitoring $\xi$, which captures unobserved factors such as hassle costs and privacy concerns. In renewal periods, it includes additive inertia costs for switching plans $\eta^{y}$ and for switching firms $\eta^{f}$.  The OOP corresponds to the portion of accident expenditures that exceeds the coverage level of plan $y$: $e({A}_{t},d) = \max(A_{t} - y_{d}, 0)$.

At the start of each period, the consumer observes their current price menu, but they do not yet know whether an accident will occur, or what the premiums in future periods will be. To account for the stochastic nature of these variables, we assume that the consumer correctly anticipates their distributions and can take expectations with respect to them. As we detail in \Cref{sec:econometric-model}, OOP is drawn from a Poisson-Pareto distribution that is truncated below by each plan's coverage limit.

While we observe the price menus offered to consumers across firms and plans, our estimation requires pricing \textit{rules}---that is, a model of how price menus would evolve at renewals based on the (potentially different) realization of stochastic accidents and monitoring scores.
Let $f^0$ denote the monitoring firm; the renewal price menu facing its customers is given by:
\begin{equation}\label{eq:renewal-rate-overview}
\begin{split}
    p_{t+1}(d | \tilde{x}_{t}, A_{t}, s, \nu_t) = &
    \begin{cases}
    p_t(f^0,y_d | \tilde{x}) \cdot (1+r_c)^{\mathbf{1}\{A_t > 0\}} \cdot r_t(m_d | \tilde{x}, s, \nu_t) & \text{ if }f_d=f^0\\
    p_0(f_d,y_d | \tilde{x}) \cdot (1+r_c)^{\mathbf{1}\{A_t > 0\}} & \text{ if }f_d\not=f^0.
    \end{cases}
\end{split}
\end{equation}
Here, $\tilde{x}$ represents consumer observables that don't relate to accidents, monitoring or pricing and $r_c$ represents a fixed rate-increase factor (an \textit{accident surcharge}) for each new accident accrued.\footnote{Note that the realizations of $A_t, s$ and $\nu_t$ are observable to the firm at the time that renewal prices are offered.} As accident records are public, this affects pricing from all firms in the market. As a simplifying assumption, we assume that consumers anticipate at most one accident per period.\footnote{This simplification is motivated empirically as we discuss below. However, our model is able to allow for multiple accidents per period in principle. In this case, $r_c$ would be applied for each accident and the total OOP for the period would be a sum of $\max(A_{tn} - y_{d}, 0)$ across $n$ accidents.} Apart from the accident surcharge, renewing consumers' premium evolves based on a {renewal factor} $r_t(m | x_t, s, \nu_t)$ that depends on their observables, an additive price shock $\nu_t$, and, if they had opted into the monitoring program, their monitoring score $s$. Non-renewing consumers instead face a new-customer ($t=0$) price menu from their chosen competitor.\footnote{Consumers have a three-period planning horizon in our model, and so they need to anticipate one more price evolution at the competitor firm when considering switching to another firm after the first period. Here, instead of estimating a separate pricing rule for each competitor, we assume that the expected price evolution at competitors follows the non-monitoring renewal factor $r_t(m=0|x_t,s=0,\nu_t)$.}

Putting these pieces together, we summarize the consumer's dynamic decision problem in each period $t$. Denoting the tuple of stochastic variables $Q_t = (A_t, v_t, s)$, the consumer chooses their best option $d$ based on the expectation with respect to $Q_t$ of their consumption utility in time $t$ and the continuation value $V(d | Q_t; x_t, p_t, d_{t-1} )$, plus an idiosyncratic Type 1 extreme value shock $\epsilon_{dt}$:
%
\begin{align}\label{eq:dynamic_decision_model}
d_{t} =  \arg\max_{d}\Big\{\mathbb{E}_{Q_{t}}\Big[u(h(d|Q_t; x_t, p_t, d_{t-1})) + \delta V(d|Q_t; x_t, p_t, d_{t-1})\Big] + \epsilon_{dt}\Big\}.
\end{align}

To define the continuation value, we make five simplifying assumptions. (i)  We assume that consumers make their choices each period with respect to a finite horizon, using two-period look ahead (e.g., internalizing three periods of utility at each decision point). (ii) We assume that consumers anticipate at most one accident occurrence within a decision horizon (e.g., across three periods). (iii) We assume that after the first renewal period, consumers anticipate a renewal rate change only if an accident occurs. (iv) We assume that consumers anticipate making a decision to switch firms or plans only upon a renewal price change (i.e., at the first renewal, or upon an accident). (v) We assume that consumers' outside options consist of plans from the firm's top two competitors.

These assumptions are empirically motivated.
Fewer than 0.5\% of drivers in our sample incur multiple accidents within a range of three periods. Only 1.3\% of drivers make multiple switches---either to a new plan or a new firm---across up to 9 renewals per driver. Among observed switches, nearly half (44.2\%)  occur at the first renewal, when the firm transitions from its new-customer pricing algorithm to its renewal-pricing algorithm. The lower frequency of switching after the first renewal corresponds greater pricing stability: absent accidents, the average renewal rate change is -0.14\%.
Our choice
of a three-period planning horizon matches the average tenure of drivers in our estimation panel (2.98 periods). We discuss this in detail and conduct robustness exercises with alternative planning horizons in Sections \ref{subsec:estimation_results} and \ref{subsec:model_sensitivity}.

Finally, the extent to which consumers attrit in renewal periods may depend on the prices they expect to receive from competitors.
Because the auto insurance market is highly fragmented with over fifty firms, we focus on the top three by market share. As we show in Table \ref{tab:sum_stat}b, this includes the monitoring firm and Competitors 1 and 2, which represent the median and the lowest price-setters, respectively, among the top-five competitors for which pricing data are available. Together, these firms account for 44.4\% of the direct (non-agency) market (42.1\% of the total market), of which 11.8\% corresponds to the monitoring firm.
Given their prominence and the representativeness of their pricing, we use these two firms to approximate consumers' outside options in our demand model.

\subsection{Econometric Specification} \label{sec:econometric-model}
Our model involves several sets of parameters, each relating to a different component of drivers' expectations when making a decision: costs in the form of claims risk and renewal prices and preferences captured by utility parameters such as risk aversion and switching frictions.
We assume that consumers have rational expectations regarding the costs involved with each insurance plan under their consideration. Their utility parameters can thus be interpreted as explaining residual variation in choices beyond what can be explained by underlying costs. Because our goal is to fit micro-data choice variation as opposed to aggregate market shares, our model allows for substantial heterogeneity in each of its key parameters. Almost all latent parameters are allowed to vary with a vector of 28 observable characteristics. In addition, we model asymmetric information between drivers and insurers by endowing each driver with a \textit{private} risk type captured by a time-invariant component of accident risk that is known to the consumer but unobserved by the firm. The firm knows only the distribution of private types when setting prices, while each consumer internalizes their private type when selecting plans and the monitoring participation.

In the remainder of \Cref{sec:econometric-model}, we outline our econometric specification. We start by modeling risk: claims, OOP, a driver-specific private risk component, and a reverse moral hazard effect due to monitoring. We then move on to modeling how the monitoring score captures risk and affects renewal price. Finally, we specify utility as well as consumers' risk aversion and choice frictions. Sections \ref{sec:identification} and \ref{sec:model_estimation} provide an overview of identification and estimation.


\paragraph{Claims risk and OOP}
Following actuarial convention, we model accident arrivals as the sum of two independent Poisson draws: major ($\mathrm{maj}$) and minor ($\mathrm{min}$) ones.
\begin{equation}\label{eq-claim-occurrence-model}
\begin{aligned}
N_{it}  = N_{it}^{\mathrm{maj}} + N_{it}^{\mathrm{min}}
\quad\text{ and }\quad
\lambda_{it} = \lambda_{it}^{\mathrm{maj}} + \lambda_{it}^{\mathrm{min}},
\quad
\text{ where } N^{\bullet}_{it} \sim \operatorname{Poisson}(\lambda^{\bullet}_{it}).
\end{aligned}
\end{equation}
The severity of major accidents is fat-tailed and follows a Pareto distribution with a position---and hence a minimum---of $\$10,000$. The severity of minor accidents follows a log-normal distribution and and is typically below $\$10,000$.
\begin{align}\label{eq-claim-severity-model}
    S_{it}^{\mathrm{maj}}  \sim \text{Pareto}(10^4, \alpha^{\mathrm{maj}}_{it}) \quad \text{and} \quad \log(S^{\mathrm{min}}_{it}) \sim \mathcal{N}(\mu^{\mathrm{min}}_{it}, \sigma^{\mathrm{min}}_{it}).
\end{align}
This distinction is conceptually important because the minimum coverage option in our Illinois panel is \$40,000. As such, while any accident would trigger a rate increase for renewal prices, only major accidents can lead to OOP. Under our model, the expected value of a consumer's OOP is the product of their expected number of major accidents, as determined by $\lambda^{\mathrm{maj}}$, and the expected cost of a major accident, as determined by $\alpha^{\mathrm{maj}}$ and their chosen liability coverage. We model the severity shape parameters as generalized linear models of observables, including zip code income and vehicle class:
\begin{align}\label{eq-claim-severity-hyper-params}
     \alpha_{it}^{\mathrm{maj}} & = f(\theta^{\alpha_\text{maj}}_0 + \mathbf{x}_{it}^\prime \cdot \mathbf{\theta}^{\alpha_\text{maj}}_{1}) \\
     \mu_{it}^{\mathrm{min}} & =  \theta^{\mu_\text{min}}_0 + \mathbf{x}_{it}^\prime \cdot \mathbf{\theta}^{\mu_\text{min}}_{1}.
\end{align}
Since $\alpha^{\mathrm{maj}}<1$ implies infinite Pareto variance, we parametrize it with a scaled inverse-logit GLM where $f(x) = 1 + Logit^{-1}(x) \cdot 10$ in the estimation procedure. This bounds the domain in a numerically efficient way while allowing for a sufficiently wide range of possible values. We further assume that consumers do not anticipate accident losses that are beyond half a million dollars, which is the largest claim we observe. Taken together, the accident loss within a period $A_{it}$ is given by $N_{it}^{\mathrm{maj}}\cdot S_{it}^{\mathrm{maj}} + N_{it}^{\mathrm{min}} \cdot S_{it}^{\mathrm{min}}$, which is then split into a covered amount based on the consumer's chosen coverage level and a residual incurred by the consumer as an OOP.

\paragraph{Private risk and moral hazard}
Our analysis in \Cref{subsec:rf_selection} suggests that accident risk may change while the driver undergoes monitoring, reflecting a moral hazard effect. It also contains a persistent private component that cannot be predicted by observables but can be partially revealed through monitoring scores. Motivated by these findings, we model drivers' accident rates as follows:
\begin{align}
    \log \lambda_{it}(m,z) & = \underbrace{\hat{\lambda}(x_{it})}_\text{public base rate} + \underbrace{m\cdot\theta^{\text{mh}}_0 + z\cdot\theta^{\text{mh}}_1}_\text{moral hazard effect} + \underbrace{\epsilon^{\lambda}_{i}}_\text{private type} \label{eq:lambda_model} \; \quad \; \text{and} \\
    \text{where } \;  \hat{\lambda}(x_{it}) & = \mathbf{\theta}^{\lambda}_{0} + \mathbf{x}_{it}^\prime \cdot \mathbf{\theta}^{\lambda}_{1} \\
    \log \lambda_{it}^{\mathrm{maj}}(m,z) & = \log \lambda_{it}(m, z) + {\mathbf{\theta}^{\mathrm{maj}}_{0}}, \label{eq:lambda_model2}
\end{align}
The accident rate for consumer $i$ in period $t$ is the sum of two components: a public signal modeled as a linear combination of observables $\mathbf{x}_{it}$, and a private error term $\epsilon^{\lambda}_{i}$ that is consumer-specific, time-persistent, and unobservable to firms. When consumers opt into monitoring ($m=1$) in the first period, we model the moral hazard effect as a mechanical shift in their accident rate $\lambda$ that depends on their monitoring intensity $z$.
\Cref{eq:lambda_model2} models a consumer's major accident rate as a scalar fraction ($\exp(\theta^{\mathrm{maj}}_0)$) of their total accident rate.\footnote{As major accidents are very rare, we do not model heterogeneity in the fraction of major accidents.}


\paragraph{Monitoring Score} If a consumer opts into monitoring, their driving is tracked for several months, yielding a large set of unstructured data that the firm synthesizes into a one-dimensional monitoring score through a predictive algorithm. As we showed in \Cref{subsec:rf_selection}, the monitoring score is positively correlated with private risk. However, there remains a substantial amount of random variation due to noise in the raw driving data and the interpretation of the data through the algorithm. To capture both of these forces, we model the monitoring score of an individual $i$ as follows:
\begin{align}\label{eq:score-def-main}
     s_{i}(m, M) = \theta^{s}_0 + {\theta}^{s}_{1} \cdot \epsilon^{\lambda}_{i} + {\theta}^{s}_{2} \cdot \hat{\lambda}_{it=0} + \epsilon^{s}_{i} \quad \text{ where }\epsilon^{s}_{i} \sim N(0, \sigma^{s}).
\end{align}
The monitoring score linearly interpolates between a driver's public risk type $\hat{\lambda}_{it=0}$ (based on observables) and their private type $\epsilon^{\lambda}_{i}$, plus an idiosyncratic mean-zero shock that captures the noise in the monitoring process (and thus, the added reclassification risk). The relative magnitudes of the coefficients on the risk components determine the extent to which the monitoring score updates the firm's prior beliefs about a driver's risk type. As such, a privately safer driver would expect their score to reflect a lower risk relative an average driver with the same observables. However, the idiosyncratic shock $\epsilon^{s}_{i}$, which accounts for noise in the monitoring process, introduces additional reclassification risk, which may prevent some risk-averse drivers from opting in even if they expect a discount on average and suffer no monitoring disutility.


\paragraph{Renewal Prices}
The renewal factors, $r_t(m|x_{it},s_i,\nu_{it})$ as described in Equation \ref{eq:renewal-rate-overview}, govern how consumers' price menus would have evolved under different monitoring and claim realizations. We model them as follows:
\begin{align}\label{eq:renewal-factor-score}
r_t(m|x_{it}, s_i, \nu_{it}) & =
\begin{cases}
\theta^{r0}_0 + \theta^{r0}_1 \cdot RC_{it} + \nu_{it}\cdot \sigma^{r0}, & \text{if } t = 0\text{ and }m = 0, \\
\theta^{r1}_0 + \theta^{r1}_1 \cdot RC_{it} + \theta^{r1}_2 \cdot s_{i} + \nu_{it}\cdot \sigma^{r1}, & \text{if } t = 0\text{ and }m = 1, \\
\theta^{r}_0 + \theta^{r}_1 \cdot RC_{it} + \nu_{it}\cdot \sigma^{r}, & \text{if } t > 0, \\
\end{cases}\\
\text{where } \nu_{it} &\sim N(0, 1). \nonumber
\end{align}

A consumer $i$ who faces a price $p$ for a given plan in period $t$ expects to face a renewal price of $p \cdot r_t(m|x_{it}, s_i, \nu_{it})$ for period $t+1$ if they do not incur an accident. The first renewal factor ($r_{t=0}$) is distinct from that of subsequent periods in two ways. First, it is generally higher---in large part because it corresponds to a shift from the firm's initial-period pricing algorithm to a distinct renewal pricing algorithm. Second, it incorporates additional information about private risk obtained through the monitoring program. After the first renewal ($t>0)$, monitoring no longer plays an active role in pricing and the renewal factors $r_{t>0}$ tend to vary much less. Nonetheless, as renewal factors enter prices multiplicatively, monitored drivers carry forward their discount or surcharge into perpetuity.

To capture these different renewal factor paths in a parsimonious manner, we model $r_t$ as a linear function of a consumer's risk class $RC$---the actuarially fair premium based on the firm's risk rating, which condenses the information contained in observables $x$ for the sake of pricing and is taken as data---as well as an intercept, their monitoring opt-in choice $m$ and score $s$ (if relevant), and an idiosyncratic mean-zero error term $\nu$. We estimate a different set of coefficients $\theta^{r}$ and error variance $\sigma^{r}$ for the first renewal with and without monitoring, and for subsequent renewals, as characterized in Eq. \ref{eq:renewal-factor-score}.

Taken together, we have six parameters to capture each pricing regime and four to capture each monitoring regime. We observe three pricing regimes and three monitoring regimes over our research period, captured by a total of thirty parameters.

\paragraph{Risk Aversion and Utility Specification} We evaluate utility using a normalized second-order Taylor approximation of von Neumann-Morgenstern utility around income $w$. This approach is similar to \textcite{Cohen2007}, and belongs to a broader class of ``negligible third-derivative'' models \parencite{barseghyan2018estimating}. The main advantage is that it simplifies the nonlinearity of the utility function, allowing us to numerically integrate over stochastic accident occurrence states while analytically reducing the stochasticity associated with claim severity, monitoring score, and renewal prices.
\begin{equation}
    u_{it}(h) = (w_{it} + h) - \frac{\gamma_{it}}{2} (w_{it} + h)^2.
\end{equation}
%
We use drivers' home zip code income in the corresponding calendar year as a proxy for $w$. This is supported by many recent surveys.\footnote{The median U.S. household only has around \$1,700 in liquid wealth \parencite{kaplan2014wealthy}. We halve the annual zipcode income to match the six-month period.} Moreover, given our utility specification, any heterogeneity in income levels can be re-parameterized to be captured by risk aversion $\gamma_{it}$, which we allow to vary flexibly based on a rich set of observables \parencite{jeziorski2014adverse}:
\begin{equation}
    \gamma_{it} = \exp( \theta^{\gamma}_0 + \mathbf{x}_{it}^\prime \cdot \mathbf{\theta}^{\gamma}_{1}).
\end{equation}
%
\paragraph{Choice Frictions} Our demand model allows for four types of utility frictions: (i) a fixed effect for the default plan featuring the mandatory minimum coverage; (ii) a monitoring disutility; (iii) a plan-switching cost; and (iv) a firm-switching cost.  In each case, the friction enters consumers' utility as an additive scalar parameter that varies with consumer observables and applies for a subset of menu options. In this way, these frictions explain the mean residual variation in choices that are not explained by the direct plan values.
To capture this flexibly across different types of drivers, we model each friction as a linear regression of observables.

\subsection{Identification\label{sec:identification}}

We now provide an informal discussion of the variation in our data that allows us to identify the parameters of our model.

\paragraph{Claims Risk Parameters} We observe a panel of claim events (and non-events) for a large number of consumers. Under a Poisson model, each individual's average claims rate is a sufficient statistic for their risk parameter. Because accidents are rare, however, we gain precision by pooling in the cross-section through the regression function described in \Cref{eq:lambda_model}. This cross-sectional variation is especially important for identifying the severity parameters (based on observed claim severity), as many drivers do not experience a major claim in our sample. The private risk component $\epsilon^{\lambda}_{i}$ can be thought of as a random intercept in the Poisson regression function. In this sense, our panel of claims is sufficient to identify it in distribution. In addition, for consumers who entered monitoring, we also observe monitoring scores that are informative of private risk as described in \Cref{eq:score-def-main}.

\paragraph{Utility Parameters}
The identification of utility parameters such as consumers' risk aversion and utility frictions relies on rich variation in the pricing and offerings of insurance coverage and monitoring.
For each consumer, we observe a menu of choices with different characteristics and prices determined by the timing of their arrival at the firm. Our claims and price models map these choices into expected costs for the current and subsequent periods. The utility parameters then rationalize the cost trade-offs that induce choice shares that match those observed in the data.

Our Illinois panel covers several baseline pricing regimes and telematics pricing regimes. As we demonstrate in \Cref{tab:app-est-pricescore} and Figures \ref{fig:score-vs-risk-by-regime} and \ref{fig:R-vs-risk-by-regime}, each regime exhibits a distinct relationship between claims risk and renewal pricing. These relationships arise from changes in both the assignment and the pricing of drivers' risk classes and monitoring scores as part of major rate revisions approved by state commissioners and updates to the monitoring algorithms.
As we showed in \Cref{subsec:rf_demand}, the average price elasticity estimated within narrow windows around these revision events---during which consumers' exposure to the new rates is plausibly random---is similar to the average elasticity across the full sample (\Cref{eq:rf_demand_1S} and \Cref{fig:rf-demand}). As such, although our structural analysis uses the whole sample in order to leverage other sources of variation, we expect that restricting attention to these quasi-experimental windows would yield similar results as well.

We also observe variation in the menu of plans available to different consumers and over time.
For instance, consumers who arrived before monitoring was introduced or whose vehicles were older than 1995 were not eligible for the monitoring program.
Moreover, a change to the mandatory minimum requirement was implemented during our sample period, exogenously changing the baseline coverage from \$40,000 to \$50,000.

The identification of the risk aversion $\gamma_{it}$ and friction parameters closely follows the literature. Different $\gamma_{it}$ values imply different gradients of $\Delta u_{it}$ across the multiple coverage options observed in the data. Conditional on risk parameters, risk aversion is therefore identified from how coverage shares vary with the contract space and with pricing gradients across coverage options, as studied in \textcite{kircher2015inferring, barseghyan2018estimating}. In addition, similar to \textcite{Handel2015}, consumers' risk aversion affects their plan choice both because higher coverage reduces the risk of OOP (a static channel) and because more price fluctuations create reclassification risk (a dynamic channel). The friction parameters---plan-specific fixed effects and path-dependent switching costs---are modeled analogously to \textcite{Handel2013}, which rationalizes features of the data not captured by risk and risk aversion alone, including the observed share of the mandatory minimum plan, the attrition and plan-switching rates in renewals, and the monitoring opt-in rate.

\paragraph{Claim Severity and Renewal Pricing Parameters} Our models of claim severity and renewal pricing are meant to parsimoniously represent the ways in which these empirical objects co-vary with observables in our data. Per actuarial standards, claim severity is considered to be orthogonal to accident risk, and largely depends on the characteristics of the vehicle being driven and the area being driven in. As such, it can be estimated separately from our claim occurrence and demand models. We accomplish this through a standard generalized linear model of drivers' observable characteristics that is identified in the cross-section. Since major claims are sometimes capped above by coverage choice, we incorporate such truncation, whenever observed, in the Pareto maximum likelihood estimation.

Similarly, our renewal pricing parameters model the pricing rule based on observed renewal price changes. We fit this using straightforward OLS regressions based on the realized price menu offered to each consumer throughout their tenure with the firm.  This is a fixed algorithm based on observable characteristics in practice, and so the parameters are identified in the cross-section.


\subsection{Estimation\label{sec:model_estimation}}
We estimate our model of insurance cost and demand using a penalized maximum likelihood procedure with a Bayesian bootstrap in two steps. In the first step, we estimate our model of claim severity conditional on an accident as characterized in \Cref{eq-claim-severity-model} and our model of renewal prices as characterized in \Cref{eq:renewal-factor-score}. In the second step, we take the parameter estimates from the claim severity and renewal price models as inputs, and jointly estimate our model of claim occurrence, monitoring score and insurance demand. Our choice to estimate claim severity and renewal prices separately from demand reflects a key assumption: consumers have rational expectations about OOP and the evolution of prices. That is, rather than allow the distributions of severity and renewal prices to be influenced by consumers' revealed preferences (and beliefs) through their insurance plan choices, we treat these objects as objective facts and estimate them from the realized claims and realized renewal price panels, respectively, alone.

By contrast, our joint estimation of claim occurrence, monitoring score and insurance choice ensures that drivers' private information regarding their risk is consistent across the different parts of the model, such that (for instance) a driver's decision to enter monitoring is informative of their underlying risk type. The contributions of claim occurrences to the likelihood function are given directly by Poisson distribution as characterized in \Cref{eq-claim-occurrence-model} and the contributions of the monitoring scores are given by the Gaussian error distribution characterized in \Cref{eq:score-def-main}. Because claim occurrences are rare, we complement our individual likelihood contributions with the likelihood of aggregate moments based on the distributions of the population averages of claim occurrence (major and minor) and choice (coverage, monitoring, and attrition) shares in our sample.\footnote{Augmenting our maximum likelihood objective with population moment likelihoods is similar to the regularization procedures described in \textcite{chernozhukov2003mcmc, beaulac2023moment} and is meant to prevent overfitting the likelihood of sparse claims observations on the basis of the dense insurance choice observations.} The likelihood contribution of each observed insurance choice $d_{it}$ is given by the mixed logit choice probability of $d_{it}$ induced by the Type 1 extreme value shock in the dynamic utility specification characterized in \Cref{eq:dynamic_decision_model}.

Although in principle, it is possible to compute standard errors directly from the Maximum Likelihood estimator, the multi-step estimation procedure and high dimensionality of our parameter space make this approach computationally intractable. Instead, we compute standard errors through a joint Bayesian bootstrap procedure. For each of our estimation components (e.g. the claim severity estimator, the renewal price estimator and the joint claim risk, monitoring score and insurance choice estimator) we draw independent observation weights for each observation in our panel according to a Dirichlet distribution 100 times. We then perform the full estimation procedure for each of the 100 sets of draws and compute standard errors for each parameter based on the standard deviation among draw estimates.
